\documentclass[11pt]{article}

\usepackage[T1]{fontenc}
\usepackage[utf8]{inputenc}
\usepackage{lmodern}
\usepackage{amsmath,amssymb,amsthm}
\usepackage{geometry}
\usepackage{booktabs}
\usepackage{array}
\usepackage{enumitem}
\usepackage{xcolor}
\usepackage[colorlinks=true,citecolor=blue,linkcolor=blue,urlcolor=blue]{hyperref}
\usepackage{microtype}

\geometry{margin=1in}
\urlstyle{same}
\sloppy

\newcommand{\Lean}{\textsc{Lean}}
\newcommand{\Mathlib}{\textsc{Mathlib}}
\newcommand{\code}[1]{\texttt{\detokenize{#1}}}
\newcommand{\codestack}[2]{\shortstack[l]{\code{#1}\\\code{#2}}}
\newcommand{\foundations}{[\code{propext}, \code{Classical.choice}, \code{Quot.sound}]}

\title{A Lean companion to Classical Lottery in Action}
\author{Jingyuan Li\thanks{Lingnan University, Hong Kong. jingyuanli@ln.edu.hk},
\ Ilia Tsetlin\thanks{INSEAD, Singapore. ilia.tsetlin@insead.edu},
\ Fan Wang\thanks{ESSEC Business School, Singapore. wang.fan@essec.edu}}
\date{May 22, 2026}

\begin{document}
\maketitle

\begin{abstract}
This note documents the \Lean~4 artifact \code{ClassicalLotteryInAction.lean} accompanying the published Management Science paper \emph{Classical Lottery in Action: Quantifying Risk and Evaluating Uncertainty}.  The artifact formalizes the local classical-lottery infrastructure, the average-utility consequences used for classical lotteries, matching-frequency constructions for two-prize acts, and the public wrapper statements for the smooth-ambiguity representation and ambiguity-attitude results.  Its main design choice is audit transparency: non-local representation-theoretic ingredients are exposed as named bridge hypotheses, while the public theorem path is checked by \Lean.  The axiom audit for the five public results reported here contains only \foundations; the wrapper hypotheses remain ordinary theorem assumptions, not hidden axioms.  A companion spin-out, \code{MechanizedDecisionTheoryWakkerDK.pdf}, is cited as ongoing work on the Wakker/Debreu--Koopmans representation layer imported by the Management Science artifact.
\end{abstract}

\section{Purpose of the note}

The economic and conceptual contribution is the published Management Science article \cite{LiTsetlinWang2026MS}.  This short companion note has a narrower purpose: it explains what the \Lean artifact proves, which theorem-shaped ingredients it deliberately leaves as named bridges, and how to reproduce the artifact-level axiom audit.

The source file is
\[
\code{ClassicalLotteryInAction.lean}.
\]
It formalizes classical lotteries as non-empty lists, acts as state-indexed lotteries, and preferences as binary relations on acts.  The file then develops the operations and predicates used in the paper: concatenation, replication, deletion of matched sub-lotteries, relative frequencies, average utility, two-prize matching frequencies, smooth representations, and ambiguity-attitude predicates.

The artifact is not a monolithic formal proof of every representation theorem invoked in the paper.  Instead, it follows a named-bridge architecture: local algebraic, order-theoretic, and finite-sum facts are proved directly, while monograph-scale or non-local representation steps are imported through explicit \code{Prop}-valued hypotheses.  This makes the formal boundary inspectable.

\section{Artifact surface}

The public-facing Lean statements tracked by this note are as follows.

\begin{itemize}[leftmargin=*]
\item \code{prop_average_utility}: Proposition~1, the average-utility characterization for classical lotteries.  The easy direction from representation to Cancellation and Archimedeanity is proved; the representation-theorem direction is the named bridge \code{AverageUtilityHardDirection}.
\item \code{lem_gap_filling}: the two-prize matching-frequency gap-filling lemma.  Its reverse implication uses the named strict-separation bridge \code{MatchingFrequencyStrictSeparation}.
\item \code{thm_smooth_model}: the smooth-representation theorem under explicit Wakker/Debreu--Koopmans-style sufficiency and regularity wrappers.
\item \code{matching_freq_smooth_formula}: the normalized two-prize formula for behavioral matching frequency, with the inverse-\(\psi\) supremum step named as \code{MatchingFrequencySmoothFormulaBridge}.
\item \code{prop_aversion_or_seeking}: the ambiguity-aversion/ambiguity-seeking curvature statement.  The easy directions from curvature to behavior are proved; behavioral recovery of curvature is isolated as \code{AmbiguityAttitudeCurvatureBridge}.
\end{itemize}

The source also defines Axiom~8, \code{Solvability}, and proves boundary lemmas about solvability, but the five rows above are the main theorem-level audit targets corresponding to the formalization-status paragraph in the published paper.

\section{Named-bridge interface}
\label{sec:bridges}

Table~\ref{tab:bridges} records the bridge interface exposed by \code{ClassicalLotteryInAction.lean}.  The important point is that these are Lean assumptions at theorem boundaries, not additional Lean axioms.

\begin{table}[ht]
\centering
\small
\begin{tabular}{@{}>{\raggedright\arraybackslash}p{0.27\textwidth}>{\raggedright\arraybackslash}p{0.32\textwidth}>{\raggedright\arraybackslash}p{0.33\textwidth}@{}}
\toprule
Bridge & Lean role & Mathematical content named by the bridge \\
\midrule
\shortstack[l]{\code{AverageUtility}\\\code{HardDirection}\\\code{P}} & Input to \code{prop_average_utility} & The vNM/Wakker-style hard direction: Cancellation plus Archimedeanity produce an average-utility representation on constant classical lotteries. \\
\addlinespace
\codestack{MatchingFrequency}{StrictSeparation P} & Input to \code{lem_gap_filling} & On two-prize acts, strict preference implies strict separation of matching frequencies.  The file further decomposes it into \code{TwoPrizeDenseness} and \code{TwoPrizeStrictRelFreq}. \\
\addlinespace
\codestack{SmoothRepresentation}{Sufficiency Pref} & Forward direction of \code{thm_smooth_model} & From non-triviality and Axioms~1--7, obtain a smooth representation.  This is the Wakker/Debreu--Koopmans representation-and-patching import. \\
\addlinespace
\codestack{SmoothRepresentation}{Regularity Pref} & Reverse regularity part of \code{thm_smooth_model} & From a smooth representation, recover Denseness, Continuity, and Consistent Aggregation; Weak Order, Cancellation, Archimedeanity, and Monotonicity are derived locally. \\
\addlinespace
\shortstack[l]{\code{MatchingFrequency}\\\code{SmoothFormulaBridge}\\\code{Pref R}} & Input to \code{matching_freq_smooth_formula} & Under \(u(x)=1\), \(u(y)=0\), identify the behavioral matching frequency with \(\psi^{-1}\!\left(\int \psi(r_x(f(s)))\,dP\right)\).  The integrand simplification is proved separately. \\
\addlinespace
\shortstack[l]{\code{AmbiguityAttitude}\\\code{CurvatureBridge}\\\code{Pref R}} & Input to \code{prop_aversion_or_seeking} & Converse Debreu--Koopmans-style curvature recovery: ambiguity aversion implies concavity of \(\psi\), and ambiguity seeking implies convexity. \\
\bottomrule
\end{tabular}
\caption{Named bridges in \code{ClassicalLotteryInAction.lean}.}
\label{tab:bridges}
\end{table}

The matching-frequency bridge has an additional useful reduction inside the artifact.  \code{TwoPrizeStrictRelFreq} is theorem-backed once \code{AverageUtilityHardDirection} supplies an average-utility representation, and \code{TwoPrizeDenseness} is derived from primitive \code{Denseness} and \code{Monotonicity} under the same hard-direction bridge.  Thus the \code{_of_averageUtilityHardDirection} route records a single named path from the Proposition~1 representation spine to the gap-filling strict-separation hypothesis.

\section{Axiom audit}
\label{sec:audit}

The source file contains explicit \texttt{\#print axioms} commands for the five public theorem targets.  Running
\[
\code{lake env lean "ClassicalLotteryInAction.lean"}
\]
prints the audit in Table~\ref{tab:audit}.  The result is intentionally uniform: every row depends only on the standard Lean/mathlib foundations \foundations.

\begin{table}[ht]
\centering
\small
\begin{tabular}{@{}>{\raggedright\arraybackslash}p{0.29\textwidth}>{\raggedright\arraybackslash}p{0.39\textwidth}>{\raggedright\arraybackslash}p{0.24\textwidth}@{}}
\toprule
Public theorem & Paper role & Axiom dependencies reported by Lean \\
\midrule
\code{prop_average_utility} & Proposition~1: average utility for classical lotteries & \foundations \\
\addlinespace
\code{lem_gap_filling} & Lemma~1: matching-frequency gap filling on two-prize acts & \foundations \\
\addlinespace
\code{thm_smooth_model} & Theorem~1: smooth-representation wrapper theorem & \foundations \\
\addlinespace
\codestack{matching_freq}{_smooth_formula} & Equation/formula for normalized smooth matching frequency & \foundations \\
\addlinespace
\code{prop_aversion_or_seeking} & Proposition~2: ambiguity attitudes and curvature & \foundations \\
\bottomrule
\end{tabular}
\caption{Axiom audit from the artifact.  The bridge hypotheses in Section~\ref{sec:bridges} are theorem assumptions and therefore do not appear as global Lean axioms.}
\label{tab:audit}
\end{table}

This audit should be read together with the named-bridge table.  The statement ``no extra axioms'' does not mean that the Wakker/Debreu--Koopmans representation theorem, strict two-prize separation, inverse-\(\psi\) matching-frequency identification, or curvature-recovery arguments have been magically proved.  It means that the artifact exposes those obligations explicitly as parameters, and the checked code downstream of those parameters introduces no additional global assumptions.

\section{Relation to the Wakker/Debreu--Koopmans spin-out}

The file \code{MechanizedDecisionTheoryWakkerDK.pdf} is a separate work-in-progress note on the Wakker additive-representation and Debreu--Koopmans concavity layer \cite{MechanizedWakkerDK2026}.  Its certificate vocabulary is complementary to the bridge surface in \code{ClassicalLotteryInAction.lean}: it studies construction certificates such as \code{hConstruct}, \code{hglobal}, \code{haff}, \code{hConc}, and \code{hPairConc}, together with a certificate checklist in \code{WakkerDebreuKoopmans.lean}.  That spin-out is the appropriate home for the multi-month theorem-prover work behind the imported Wakker/Debreu--Koopmans layer.

The Management Science artifact therefore has a stable publication role: it states and checks the classical-lottery reduction, matching-frequency infrastructure, and smooth-model theorem boundaries used by the published paper.  The Wakker/Debreu--Koopmans spin-out has a different role: to progressively discharge the imported representation certificates without changing the public interface consumed here.

\section{Reproducibility notes}

The audit command used for this note was run from the repository root on Windows:
\[
\code{lake env lean "ClassicalLotteryInAction.lean"}.
\]
The relevant output was:
\begin{quote}\small
\code{'ClassicalLottery.Preference.prop_average_utility' depends on axioms: [propext, Classical.choice, Quot.sound]}\\
\code{'ClassicalLottery.Preference.lem_gap_filling' depends on axioms: [propext, Classical.choice, Quot.sound]}\\
\code{'ClassicalLottery.Preference.thm_smooth_model' depends on axioms: [propext, Classical.choice, Quot.sound]}\\
\code{'ClassicalLottery.Preference.matching_freq_smooth_formula' depends on axioms: [propext, Classical.choice, Quot.sound]}\\
\code{'ClassicalLottery.Preference.prop_aversion_or_seeking' depends on axioms: [propext, Classical.choice, Quot.sound]}.
\end{quote}

\begin{thebibliography}{9}

\bibitem{LiTsetlinWang2026MS}
Jingyuan Li, Ilia Tsetlin, and Fan Wang.
\newblock \emph{Classical Lottery in Action: Quantifying Risk and Evaluating Uncertainty}.
\newblock Management Science, 2026.

\bibitem{MechanizedWakkerDK2026}
Jingyuan Li, Ilia Tsetlin, and Fan Wang.
\newblock \emph{Certificate-Driven Mechanization of Wakker--Debreu--Koopmans Additive Representation in Lean 4}.
\newblock Source draft \code{MechanizedDecisionTheoryWakkerDK.pdf}, ongoing work, 2026.

\bibitem{Wakker1989}
Peter P. Wakker.
\newblock \emph{Additive Representations of Preferences: A New Foundation of Decision Analysis}.
\newblock Kluwer Academic Publishers, 1989.

\bibitem{DebreuKoopmans1960}
Gerard Debreu and Tjalling C. Koopmans.
\newblock Additively decomposed quasiconvex functions.
\newblock \emph{Mathematical Programming and Related Techniques}, 1960.

\end{thebibliography}

\end{document}
